\chapter{Mesenteric adenitis}

\section*{Definition and Overview}
Mesenteric adenitis, also known as mesenteric lymphadenitis, is an inflammatory condition affecting the lymph nodes situated within the mesentery—the double layer of peritoneum that attaches the intestines to the posterior abdominal wall. Typically presenting as a benign, self-limiting condition, it is a common cause of acute abdominal pain in children and adolescents. Its primary clinical significance lies in its presentation: the acute right lower quadrant abdominal pain, fever, and tenderness characteristic of mesenteric adenitis closely mimic the clinical signs of acute appendicitis. Distinguishing between these two conditions is essential to prevent unnecessary surgical interventions while ensuring that acute surgical emergencies are not missed.

\section*{Epidemiology}
Mesenteric adenitis predominantly affects paediatric and adolescent populations, most commonly children under the age of 15 years, although it can occasionally occur in young adults. It exhibits a slight male predominance. Paediatric clinical audits indicate that mesenteric adenitis is responsible for approximately 15\% to 20\% of cases where children present to emergency departments with acute right lower quadrant abdominal pain suspicious of appendicitis. The condition shows a seasonal pattern, with a higher incidence during late autumn, winter, and early spring, correlating with the peak seasonal circulation of viral respiratory and gastrointestinal tract infections in the community.

\sectionsection*{Aetiology and Pathophysiology}
The inflammation of the mesenteric lymph nodes is typically a reactive response to a primary infective process elsewhere in the body.
\begin{itemize}
    \item \textbf{Infectious triggers:} The most frequent triggers are viral infections of the upper respiratory tract or gastrointestinal system. Common viral agents include \textbf{adenovirus}, \textbf{enterovirus}, \textbf{rotavirus}, and \textbf{Epstein-Barr virus (EBV)}. Bacterial infections can also cause the condition, with \textbf{\textit{Yersinia enterocolitica}} being a classic foodborne bacterial pathogen associated with severe mesenteric lymphadenitis. Other bacterial agents include \textbf{\textit{Campylobacter}} species, \textbf{\textit{Salmonella}} species, and \textbf{\textit{Shigella}} species.
    \item \textbf{Pathology:} Pathogens transit from the intestinal mucosa via the lymphatic vessels to the regional mesenteric lymph nodes, primarily in the ileocecal region. This leads to lymphoid hyperplasia, vascular congestion, and nodal enlargement.
    \item \textbf{Anatomical localization:} The high concentration of lymph nodes around the terminal ileum and cecum explains why the inflammatory response and subsequent abdominal tenderness are localised predominantly to the right lower quadrant of the abdomen.
\end{itemize}

\section*{Clinical Features}
The symptoms of mesenteric adenitis are variable but characteristically mimic acute intra-abdominal pathology:
\begin{itemize}
    \item \textbf{Abdominal pain:} Typically localized in the right lower quadrant or periumbilical region. The pain is often described as dull or crampy and, unlike appendicitis, it may fluctuate in intensity and shift slightly when the patient moves.
    \text{Fever:} Low-to-moderate grade fever (typically between 38$^\circ$C and 39$^\circ$C), which often precedes the onset of abdominal pain.
    \item \textbf{Gastrointestinal symptoms:} Mild nausea, occasional vomiting, anorexia, and mild diarrhoea may accompany the pain.
    \item \textbf{Antecedent infection:} A history of an upper respiratory tract infection (sore throat, rhinorrhoea, earache) or gastroenteritis in the preceding one to two weeks.
    \item \textbf{Generalised lymphadenopathy:} Swollen, mildly tender lymph nodes in the cervical region (neck) may be noted on examination, reflecting a systemic immune response.
\end{itemize}

\section*{Differential Diagnosis}
Because of the risk of misdiagnosis, mesenteric adenitis must be carefully differentiated from several conditions:
\begin{itemize}
    \item \textbf{Acute appendicitis:} The most critical differential diagnosis. Appendicitis pain typically starts periumbilically, migrates to the right lower quadrant, becomes progressively constant and severe, and is associated with localized guarding and rebound tenderness.
    \item \textbf{Intussusception:} A medical emergency in infants (typically 3 to 24 months old) presenting with severe, episodic colicky pain, vomiting, and classically, ``red currant jelly'' stools.
    \item \textbf{Inflammatory bowel disease (IBD):} Crohn's disease can present with terminal ileitis, right lower quadrant pain, and mesenteric lymphadenopathy, but it is characterized by a chronic history of weight loss, recurrent diarrhoea, and growth failure.
    \item \textbf{Urinary tract infection (UTI):} Can cause lower abdominal pain and fever in children, ruled out via urinalysis.
    \item \textbf{Gynaecological pathologies:} In adolescent females, conditions such as pelvic inflammatory disease (PID), ruptured ovarian cysts, or ovarian torsion must be considered.
\end{itemize}

\section*{When to Seek Medical Care}
Parents should seek medical advice if their child develops new or worsening abdominal pain, persistent fever, or is unable to tolerate fluids.

\textbf{Call 000 (or local emergency services) for:}
\begin{itemize}
    \item Severe, constant, or worsening abdominal pain where the child cannot stand straight, or if the abdomen is rigid, hard, or extremely tender to touch.
    \item Extreme lethargy, difficulty waking, or confusion.
    \item Signs of severe dehydration, such as crying without tears, no wet nappies for over 8 hours, sunken eyes, or cold extremities.
    \item High fever accompanied by a stiff neck, sensitivity to light (photophobia), or a dark purple, non-blanching spotty rash.
\end{itemize}

\section*{Diagnosis and Investigations}
Diagnosis is made through a combination of clinical assessment and targeted diagnostic imaging to exclude surgical emergencies.
\begin{itemize}
    \item \textbf{Clinical examination:} Palpation of the abdomen typically reveals mild, diffuse tenderness in the right lower quadrant, but without signs of peritoneal irritation (no guarding, rigidity, or rebound tenderness).
    \item \textbf{Abdominal ultrasound:} The imaging modality of choice. Diagnostic criteria for mesenteric adenitis on ultrasound include the presence of three or more enlarged mesenteric lymph nodes (measuring 8 millimetres or more in their short-axis diameter) in the right lower quadrant, in the presence of a normal-looking, compressible appendix.
    \item \textbf{Full blood count (FBC):} May show a normal or mildly elevated white cell count (leukocytosis) with lymphocytosis, which is common in viral infections.
    \item \textbf{Inflammatory markers:} C-reactive protein (CRP) or erythrocyte sedimentation rate (ESR) may be normal or mildly elevated, but are generally lower than in acute appendicitis.
    \item \textbf{Urinalysis:} Performed in all cases of acute abdominal pain to exclude a urinary tract infection.
\end{itemize}

\section*{Management}
Because mesenteric adenitis is a self-limiting viral-reactive condition, treatment is supportive and focused on symptom relief.
\begin{itemize}
    \item \textbf{Symptomatic pain relief:} Paracetamol or ibuprofen can be administered to manage fever and alleviate abdominal discomfort. Salicylates (aspirin) must be avoided in children due to the risk of Reye's syndrome.
    \item \textbf{Adequate hydration:} Encouraging frequent, small sips of water, oral rehydration solutions, or diluted juices to prevent dehydration, especially if vomiting or diarrhoea is present.
    \text{Rest and recovery:} Allowing the child to rest and avoid strenuous physical activities or contact sports until the abdominal pain has resolved.
    \item \textbf{Avoidance of routine antibiotics:} Antibiotics are not indicated for viral mesenteric adenitis. They are reserved for confirmed bacterial infections (such as severe \textit{Yersinia} septicaemia) or in immunocompromised patients.
    \item \textbf{Home monitoring and reassessment:} Parents must be instructed to monitor the child closely. If the pain worsens, localises, or if peritoneal signs develop, the child must be reassessed immediately to rule out appendicitis.
\end{itemize}

\section*{Complications}
Complications are rare, but clinical courses can occasionally be compromised.
\begin{itemize}
    \item Dehydration due to poor oral intake from nausea and vomiting.
    \item Unnecessary appendicectomy resulting from diagnostic ambiguity.
    \text{Suppurative lymphadenitis:} A very rare complication where the lymph node undergoes necrosis and forms an abscess, typically associated with specific bacterial infections like \textit{Yersinia}.
    \item Sepsis or systemic spread of infection, which is extremely rare and confined to highly immunocompromised patients.
\end{itemize}

\section*{Prognosis and Outcomes}
The prognosis for mesenteric adenitis is excellent. It is a benign condition that resolves spontaneously in almost all cases. Symptoms generally improve within a few days and resolve completely within one to two weeks. There are no long-term health consequences or structural damage to the intestines. Recurrence is uncommon but can occur if the child contracts a new, unrelated viral infection that triggers another reactive lymphatic response.

\section*{Prevention and Self-Care}
Prevention focuses on general hygiene measures to avoid gastrointestinal and respiratory infections.
\begin{itemize}
    \item Teach children proper hand hygiene, emphasizing washing hands thoroughly with soap and water before eating and after using the toilet.
    \item Avoid close contact and sharing cups or utensils with individuals who show symptoms of respiratory or gastrointestinal illness.
    \item Ensure all food, particularly pork and poultry products, is cooked to safe internal temperatures, and avoid consuming unpasteurised dairy products to prevent \textit{Yersinia} and \textit{Campylobacter} infections.
    \item Ensure children are up to date with all recommended vaccinations, including the rotavirus vaccine.
\end{itemize}

\section*{Sources and Further Reading}
The following reputable organisations provide further information and paediatric clinical guidelines on mesenteric adenitis:
\begin{itemize}
    \item Royal Children's Hospital Melbourne. \textit{Clinical guidelines and parent information on abdominal pain and mesenteric adenitis.} https://www.rch.org.au/
    \item Healthdirect Australia. \textit{Trusted information on the symptoms, diagnosis, and management of mesenteric adenitis.} https://www.healthdirect.gov.au/
    \item Sydney Children's Hospitals Network. \textit{Fact sheets on acute abdominal pain in children.} https://www.schn.health.nsw.gov.au/
    \item Mayo Clinic. \textit{Overview of mesenteric lymphadenitis causes and diagnosis.} https://www.mayoclinic.org/
    \item NHS (UK). \textit{Patient resource guide on mesenteric adenitis symptoms and self-care.} https://www.nhs.uk/
\end{itemize}
